% Part: proof-theory
% Chapter: natural-deduction
% Section: quantifiers

\documentclass[../../../include/open-logic-section]{subfiles}

\begin{document}

\olfileid[hi]{pt}{nat}{qua}

\olsection{नियमित \usetoken{P}{proof} और प्रतिस्थापन}

\Elim{\lexists} और~\Intro{\lforall} नियमों पर लागू होने वाली ``आइगेनचर
शर्तों'' के कारण परिमाणकों के नियम कुछ जटिलताएँ उत्पन्न करते हैं। इन अनुमानों
में !!{constant}~$c$ का कभी पुनः उपयोग न हो, यह सुनिश्चित करके इनमें से
अधिकांश जटिलताओं को सँभाला जा सकता है। निम्नलिखित परिभाषा प्रस्तुत करें।

\begin{defn}
कोई प्राकृतिक निगमन !!{proof}~$\delta$ \emph{नियमित} है यदि
\begin{enumerate}
  \item कोई आइगेनचर~$a$, एक से अधिक \Elim{\lexists} या~\Intro{\lforall}
  अनुमान के आइगेनचर के रूप में प्रयुक्त न हो, और
  \item यदि $a$ किसी \Elim{\lexists} या~\Intro{\lforall} अनुमान का आइगेनचर
  है, तो वह केवल क्रमशः गौण आधार-वाक्य या आधार-वाक्य तक जाने वाले निकटस्थ
  उप-!!{proof} में आए।
\end{enumerate}
\end{defn}

\begin{lem}\ollabel{lem:subst-regular} यदि $\delta$ नियमित है,~$!C$ पर
समाप्त होता है, $c$ ऐसा !!a{constant} है जो~$\delta$ में आइगेनचर के रूप में
प्रयुक्त नहीं, और $t$ ऐसा पद है जिसमें~$\delta$ का कोई आइगेनचर नहीं, तो~$c$
की प्रत्येक उपस्थिति को~$t$ से बदलकर~$\delta$ से प्राप्त
$\Subst{\delta}{t}{c}$ एक नियमित !!{proof} है। $\Subst{\delta}{t}{c}$ का
अंतिम-!!{formula}, $\Subst{!C}{t}{c}$ है। यदि $!A^x$,~$\delta$ की खुली
मान्यता है, तो $\Subst{!A}{t}{c}^x$,~$\Subst{\delta}{t}{c}$ की खुली मान्यता
है।
\end{lem}

\begin{proof}
  $\pheight{\delta}$ पर आगमन से। यदि $\pheight{\delta} = 0$, तो वह केवल
  मान्यता~$!A^x$ से बना है। तब $\Subst{\delta}{t}{c}$,
  ~ $\Subst{!A}{t}{c}^x$ है।

  यदि $\pheight{\delta} > 0$, तो हम~$\delta$ के अंतिम अनुमान के अनुसार
  स्थितियों में भेद करते हैं। प्रतिज्ञप्ति संयोजकों के नियमों की स्थितियाँ
  सरल हैं और अभ्यास के रूप में छोड़ी जाती हैं। हम उन स्थितियों पर विचार करते
  हैं जहाँ $\delta$ किसी परिमाणक अनुमान पर समाप्त होता है।
  \begin{enumerate}
    \item यदि अंतिम अनुमान $\Intro{\lforall}$ है, तो $\delta$ इस रूप का है
    \[
    \AxiomC{}
    \RightLabel{$\delta'$}
    \DeduceC{$!A(a,c)$}
    \RightLabel{\Intro{\lforall}}
    \UnaryInfC{$\lforall[x][!A(x,c)]$}
    \DisplayProof
    \]
    आगमन परिकल्पना से $\Subst{\delta'}{t}{c}$, $!A(a,t)$ का नियमित
    !!{proof} है। चूँकि $a$,~$\delta$ का आइगेनचर है, इसलिए $a$,~$t$ में नहीं
    आता। अतः $\Subst{\delta}{t}{c}$ का अंतिम अनुमान,
    \[
    \AxiomC{}
    \RightLabel{$\Subst{\delta'}{t}{c}$}
    \DeduceC{$!A(a,t)$}
    \RightLabel{\Intro{\lforall}}
    \UnaryInfC{$\lforall[x][!A(x,t)]$}
    \DisplayProof
    \]
    सही है: चूँकि $t$ में~$a$ नहीं है, इसलिए न निष्कर्ष न कोई खुली मान्यता
    ~ $a$ को समाहित करती है।
    \item यदि अंतिम अनुमान $\Elim{\lforall}$ है, तो $\delta$ इस रूप का है
    \[
    \AxiomC{}
    \RightLabel{$\delta'$}
    \DeduceC{$\lforall[x][!A(x,c)]$}
    \RightLabel{\Elim{\lforall}}
    \UnaryInfC{$!A(s(c),c)$}
    \DisplayProof
    \]
    आगमन परिकल्पना से $\Subst{\delta'}{t}{c}$,
    $\lforall[x][!A(x,t)]$ का नियमित !!{proof} है। (निष्कर्ष का पद~$s(c)$,
    ~ $c$ को समाहित कर सकता है, पर यह आवश्यक नहीं।)
    $\Subst{\delta}{t}{c}$ का अंतिम अनुमान,
    \[
    \AxiomC{}
    \RightLabel{$\Subst{\delta'}{t}{c}$}
    \DeduceC{$\lforall[x][!A(x,t)]$}
    \RightLabel{\Elim{\lforall}}
    \UnaryInfC{$!A(s(t),t)$}
    \DisplayProof
    \]
    सही है और $!A(s(t),t) \ident \Subst{!A(s(c),c)}{t}{c}$।
    \Intro{\lexists} की स्थिति भी इसी प्रकार है।
  \item यदि अंतिम अनुमान $\Elim{\lexists}$ है, तो $\delta$ इस रूप का है
    \[
    \AxiomC{}
    \RightLabel{$\delta_1'$}
    \DeduceC{$\lexists[x][!A(x,c)]$}
    \AxiomC{$\Discharge{!A(a,c)}{x}$}
    \RightLabel{$\delta_2'$}
    \DeduceC{$!C$}
    \DischargeRule{\Elim{\lexists}}{x}
    \BinaryInfC{$!C$}
    \DisplayProof
    \]
    आगमन परिकल्पना से $\Subst{\delta_1'}{t}{c}$ और
    $\Subst{\delta_2'}{t}{c}$ क्रमशः $\lexists[x][!A(x,c)]$ और खुली
    मान्यता~$!A(a,t)$ से $!C$ के नियमित !!{proof}s हैं। चूँकि $a$,~$\delta$
    का आइगेनचर है, इसलिए $a$,~$t$ में नहीं आता। अतः
    $\Subst{\delta}{t}{c}$ का अंतिम अनुमान,
    \[
    \AxiomC{}
    \RightLabel{$\Subst{\delta'}{t}{c}$}
    \DeduceC{$!A(a,t)$}
    \RightLabel{\Intro{\lforall}}
    \UnaryInfC{$\lforall[x][!A(x,t)]$}
    \DisplayProof
    \]
    सही है: $\Subst{!A(x,t)}{a}{x} \ident \Subst{!A(a,c)}{t}{c}$, और न
    निष्कर्ष~$!C$ न कोई खुली मान्यता~$a$ को समाहित करती है।
  \end{enumerate}
\end{proof}

\begin{prop}\ollabel{prop:regular}
यदि $\delta$, \Log{N1c} या \Log{N1i} में कोई !!a{proof} है, तो उसी संरचना
(विशेषतः उसी !!{height}) वाला कोई नियमित !!{proof}~$\delta'$ है, जो
$\delta$ से केवल आइगेनचरों के प्रतिस्थापन में भिन्न है।
\end{prop}

\begin{proof}
केवल इस प्रमाण के लिए,~$\delta$ के किसी आइगेनचर अनुमान को \emph{स्वच्छ}
कहें यदि उसका आइगेनचर किसी अन्य अनुमान का आइगेनचर न हो और उस अनुमान के
निकटस्थ उप-!!{proof} के अतिरिक्त~$\delta$ में कहीं न आए; अन्यथा उसे
\emph{मलिन} कहें। मान लें $n$,~$\delta$ में मलिन आइगेनचर अनुमानों की संख्या
है। हम~$n$ पर आगमन से दिखाते हैं कि $\delta$ को आवश्यक नियमित
!!{proof}~$\delta'$ में बदला जा सकता है। यदि $n=0$, तो~$\delta$ में सभी
आइगेनचर अनुमान स्वच्छ हैं, अतः $\delta$ नियमित है और कुछ सिद्ध नहीं करना।
अन्यथा कोई सर्वोच्च आइगेनचर अनुमान, मान लें \Intro{\lforall}, चुनें। उस पर
समाप्त होने वाला~$\delta$ का उप-!!{proof}~$\delta_1$ इस रूप का है
    \[
    \AxiomC{}
    \RightLabel{$\delta_2$}
    \DeduceC{$!A(c)$}
    \RightLabel{\Intro{\lforall}}
    \UnaryInfC{$\lforall[x][!A(x)]$}
    \DisplayProof
    \]
ध्यान दें कि $c$ आइगेनचर है, इसलिए वह~$\lforall[x][!A(x)]$ या किसी खुली
मान्यता में नहीं आता। प्रमाण $\delta_2$ नियमित है, क्योंकि उसमें कोई आइगेनचर
अनुमान ही नहीं। मान लें $a$ ऐसा !!a{constant} है जो~$\delta$ में नहीं आता।
\olref{lem:subst-regular} से $\Subst{\delta_2}{a}{c}$, $!A(a)$ का नियमित
प्रमाण है। आइगेनचर शर्त से~$\delta_2$ की किसी खुली मान्यता में~$c$ नहीं आता,
अतः~$\Subst{\delta_2}{a}{c}$ की किसी खुली मान्यता में~$a$ नहीं आता। इसलिए हम
~$\Intro{\lforall}$ लागू कर सकते हैं। यदि~$\delta$ में $\delta_1$ को प्रमाण
$\delta_1'$ से बदलें,
    \[
    \AxiomC{}
    \RightLabel{$\Subst{\delta_2}{a}{c}$}
    \DeduceC{$!A(a)$}
    \RightLabel{\Intro{\lforall}}
    \UnaryInfC{$\lforall[x][!A(x)]$}
    \DisplayProof
    \]
तो हमने एक मलिन आइगेनचर अनुमान को स्वच्छ अनुमान से बदल दिया है और एकमात्र
अंतर आइगेनचर $c$ का~$a$ से प्रतिस्थापन है। फलित प्रमाण में $n-1$ मलिन
आइगेनचर अनुमान हैं और परिणाम आगमन परिकल्पना से निष्पन्न होता है।
\end{proof}

\begin{prob}
\olref{prop:regular} के प्रमाण में उस स्थिति का वर्णन कीजिए जहाँ सर्वोच्च
आइगेनचर अनुमान~\Elim{\lexists} है।
\end{prob}

हम अभी सिद्ध प्रतिज्ञप्ति का स्पष्ट रूप से आह्वान नहीं करेंगे, किंतु मानेंगे
कि विचाराधीन सभी !!{proof}s नियमित हैं---यदि वे न हों, तो आइगेनचर बदलकर
उन्हें सदा नियमित बना सकते हैं।

\Olref{lem:subst-regular} और \olref{prop:regular}, \Log{N2c}
और~\Log{N2i} के लिए भी सत्य हैं। प्रमाण अभ्यास के रूप में छोड़ते हैं।

\end{document}
